\FrontMatterChapter{Abstract}

PatchOps is a repository-aware, safety-gated, abstention-aware repair system for GitHub Actions failures. It
combines log extraction, failure classification, fault localization, repair planning, patch generation, and a
closed CI validation loop to decide whether the system should open a fix pull request, request review, rerun
the workflow, escalate, or abstain.

The thesis investigates not only whether an AI system can repair CI failures, but under what conditions it
should decline to act. Across the evaluated method families, fault localization remains the dominant
bottleneck, with Recall@3 staying below 0.45 and peaking at 0.32. The results also show that improved
localization is necessary but not sufficient for successful repair: full-repository retrieval can raise
patch-apply rate to 1.00 on the gated path without lifting CI Pass@1 above 0.00 on that same path.

PatchOps is therefore positioned not as a state-of-the-art repair engine on raw Pass@1 alone, but as a
reproducible and safety-conscious repair workflow that emphasizes honesty, calibrated decision-making, and
closed-loop validation. On the original 20-case Safety Set it records 0/20 unsafe auto-fixes while maintaining
high abstention accuracy, and on the expanded 76-case decision-gate study the learned calibrated act-versus-hold
policy reaches 0.91 act/hold accuracy versus 0.855 for the earlier hand-set threshold policy. This framing
aligns the system with the practical needs of software engineering teams that value dependable automation over
aggressive but unsafe patch generation.

\vspace{1em}
\noindent\textbf{Keywords:} \ThesisKeywords.
\clearpage
